\documentclass[10pt,a4paper]{amsart}
\usepackage{amsmath,amssymb,amsthm,graphicx}
\usepackage{verbatim}
\newenvironment{reference}{\comment}{\endcomment}
%\newenvironment{reference}{}{}
\newenvironment{slogan}{\comment}{\endcomment}
\newenvironment{history}{\comment}{\endcomment}

% For commutative diagrams we use Xy-pic
\usepackage[all]{xy}

% We use 2cell for 2-commutative diagrams.
\xyoption{2cell}
% Theorem environments.
%
\theoremstyle{plain}
\newtheorem{theorem}[subsection]{Theorem}
\newtheorem{proposition}[subsection]{Proposition}
\newtheorem{lemma}[subsection]{Lemma}

\theoremstyle{definition}
\newtheorem{definition}[subsection]{Definition}
\newtheorem{example}[subsection]{Example}
\newtheorem{exercise}[subsection]{Exercise}
\newtheorem{situation}[subsection]{Situation}

\theoremstyle{remark}
\newtheorem{remark}[subsection]{Remark}
\newtheorem{remarks}[subsection]{Remarks}

\numberwithin{equation}{subsection}

% Macros
%
\def\lim{\mathop{\mathrm{lim}}\nolimits}
\def\colim{\mathop{\mathrm{colim}}\nolimits}
\def\Spec{\mathop{\mathrm{Spec}}}
\def\Hom{\mathop{\mathrm{Hom}}\nolimits}
\def\Ext{\mathop{\mathrm{Ext}}\nolimits}
\def\SheafHom{\mathop{\mathcal{H}\!\mathit{om}}\nolimits}
\def\SheafExt{\mathop{\mathcal{E}\!\mathit{xt}}\nolimits}
\def\Sch{\mathit{Sch}}
\def\Mor{\mathop{\mathrm{Mor}}\nolimits}
\def\Ob{\mathop{\mathrm{Ob}}\nolimits}
\def\Sh{\mathop{\mathit{Sh}}\nolimits}
\def\NL{\mathop{N\!L}\nolimits}
\def\CH{\mathop{\mathrm{CH}}\nolimits}
\def\proetale{{pro\text{-}\acute{e}tale}}
\def\etale{{\acute{e}tale}}
\def\QCoh{\mathit{QCoh}}
\def\Ker{\mathop{\mathrm{Ker}}}
\def\Im{\mathop{\mathrm{Im}}}
\def\Coker{\mathop{\mathrm{Coker}}}
\def\Coim{\mathop{\mathrm{Coim}}}

% Boxtimes
%
\DeclareMathSymbol{\boxtimes}{\mathbin}{AMSa}{"02}

%
% Macros for moduli stacks/spaces
%
\def\QCohstack{\mathcal{QC}\!\mathit{oh}}
\def\Cohstack{\mathcal{C}\!\mathit{oh}}
\def\Spacesstack{\mathcal{S}\!\mathit{paces}}
\def\Quotfunctor{\mathrm{Quot}}
\def\Hilbfunctor{\mathrm{Hilb}}
\def\Curvesstack{\mathcal{C}\!\mathit{urves}}
\def\Polarizedstack{\mathcal{P}\!\mathit{olarized}}
\def\Complexesstack{\mathcal{C}\!\mathit{omplexes}}
% \Pic is the operator that assigns to X its picard group, usage \Pic(X)
% \Picardstack_{X/B} denotes the Picard stack of X over B
% \Picardfunctor_{X/B} denotes the Picard functor of X over B
\def\Pic{\mathop{\mathrm{Pic}}\nolimits}
\def\Picardstack{\mathcal{P}\!\mathit{ic}}
\def\Picardfunctor{\mathrm{Pic}}
\def\Deformationcategory{\mathcal{D}\!\mathit{ef}}
\usepackage{cleveref}
\crefname{originalreference}{original reference}{original references}
\Crefname{originalreference}{Original reference}{Original references}
\makeatletter
\newlabel{lemma-base-change-on-invariant}{{048G}{source}{}{sourcepointer-1}{}}
\newlabel{lemma-base-change-on-invariant@cref}{{[originalreference][0][]048G}{[1][0][]source}{}{}{}}
\makeatother
\crefname{originalsection}{original section}{original sections}
\Crefname{originalsection}{Original section}{Original sections}
\makeatletter
\newlabel{spaces-groupoids-section-quotient-sheaves}{{044H}{source}{}{sourcepointer-2}{}}
\newlabel{spaces-groupoids-section-quotient-sheaves@cref}{{[originalsection][0][]044H}{[1][0][]source}{}{}{}}
\makeatother
\title{Quotient sheaves of relations commute with base change}
\author{The Stacks Project Authors}
\date{stacks-project:groupoids-quotients.tex:local-snapshot-20261012; GNU FDL 1.2 or later}
\begin{document}
\maketitle
\noindent\textit{Editorial note.} All delivered mathematical text is the author's own, copied from the pinned Stacks Project chapter file groupoids-quotients.tex; no statement or proof step was rewritten, reordered or removed.\par\medskip
\noindent\textit{Editorial note.} The counted proof is the author's complete argument: The sheaf-theoretic shadow of the base change compatibility of quotient stacks.\par\medskip
\noindent\textit{Editorial note.} The label lemma-base-change-quotient-sheaf names the counted statement, and the proof environment immediately following it in the source is delivered verbatim as its proof.\par\medskip
\noindent\textit{Editorial note.} The counted statement is carried by the author's own lemma environment in the chapter source.\par\medskip
\noindent\textit{Allowed standard prerequisites (editorial).}\begin{itemize}
\item Chapter Lemma lemma-base-change-on-invariant (navigation stub, not delivered): for a pre-relation j on U with an R-invariant morphism to X, the base change along any X-prime to X gives U-prime equal to the fibre product of X-prime with U over X, R-prime equal to the fibre product with R, and an induced pre-relation j-prime over B
\item Groupoids in Spaces, Section spaces-groupoids-section-quotient-sheaves: the quotient sheaf U/R is defined as the coequalizer of the two maps from R to U in the category of sheaves of sets on the fppf site
\item Coequalizers in a comma category: a coequalizer diagram in the category of sheaves over the big site remains a coequalizer diagram after restricting to the comma category of sheaves over X
\item Exactness of pullback in a topos: for any morphism of topoi the base change functor on sheaves preserves finite colimits, in particular coequalizers
\end{itemize}\medskip
\ifdefined\hypertarget\hypertarget{sourcepointer-1}{}\fi
\noindent\textit{Source reference.} Original reference 048G: sources/stacks-groupoids-quotients/pinned-original-source.tex, source lines 298--298. Prerequisite pointer invoked at 354; it names a standard fact of the same chapter that the counted argument uses, and every step of the counted proof built on it is retained verbatim in the delivered body.\par\medskip
\ifdefined\hypertarget\hypertarget{sourcepointer-2}{}\fi
\noindent\textit{Source reference.} Original section 044H: sources/stacks-chow/tags.txt, source lines 5400--5400. Prerequisite pointer invoked at 360; it names a standard fact of a sibling chapter that the counted argument uses, and every step of the counted proof built on it is retained verbatim in the delivered body.\par\medskip
\begin{lemma}
\label{lemma-base-change-quotient-sheaf}
In the situation of Lemma \ref{lemma-base-change-on-invariant}
there is an isomorphism of sheaves
$$
U'/R' = X' \times_X U/R
$$
For the construction of quotient sheaves, see
Groupoids in Spaces, Section \ref{spaces-groupoids-section-quotient-sheaves}.
\end{lemma}

\begin{proof}
Since $U \to X$ is $R$-invariant, it is clear that the map
$U \to X$ factors through the quotient sheaf $U/R$.
Recall that by definition
$$
\xymatrix{
R \ar@<1ex>[r] \ar@<-1ex>[r] &
U \ar[r] &
U/R
}
$$
is a coequalizer diagram in the category $\Sh$ of sheaves of sets on
$(\Sch/S)_{fppf}$. In fact, this is a coequalizer diagram in the
comma category $\Sh/X$. Since the base change functor
$X' \times_X - : \Sh/X \to \Sh/X'$ is exact (true in any topos),
we conclude.
\end{proof}

\end{document}
